% No preamble. Prefix dmt-. Exact actual pivot/Gaussian drift and trace kernels.
\section{An explicit trace formula for the pivot and Gaussian drift}
\label{dmt-section}

Fix the primitive ratio, a composite ratio \(R\), the actual
composite ordinary average, its selected pivots and the direct-\(R\)
internal contour gauge. The preceding homogeneous classical
depth \(M\) is arbitrary. All finite-regulator formulas in this
section are exact. Their infinite-volume interpretation uses one
chosen joint local-limit system with the previously proved uniform
exponential bounds. No numerical value, sign, old-depth limit or
universal beta coefficient is asserted.

\subsection{The exact external factors}
\label{dmt-definitions-section}

Use real Lie coordinates orthonormal for the normalized trace,
and their bilinear complexification. Write current link angles
as \(a\) and final coarse link angles as \(w\). Retained-coordinate
matrix traces use the actual Euclidean counting coordinates of
the Gaussian integral. They carry no extra old fine-volume factor.

Let \(F_R\) be the actual composite ordinary group average. Its
relative constraint and differential are
\begin{equation}
 F_a(B)=\frac1i\log\!\left[
     F_R(e^{iB}e^{ia})F_R(e^{ia})^{-1}\right],
 \qquad C_a=D_BF_a(0).
 \label{dmt-relative-constraint}
\end{equation}
For the unique selected pivot \(p(c)\) of a final bond \(c\),
\begin{equation}
 (C_a)_{c,p(c)}=R^{1-d}\pi_c(a),\qquad
 \pi_c(0)=I.
 \label{dmt-pivot-block}
\end{equation}
The actual pivot insertion, retained insertion and chart satisfy
\begin{align}
 C_ah_a&=I,\qquad
 K_a=E_{\rm free}-h_aC_aE_{\rm free},\notag\\
 D_z\Phi_a(0)&=K_a,\qquad
 D_z^2\Phi_a(0)[u,v]
       =-h_aD_B^2F_a(0)[K_au,K_av].
 \label{dmt-chart-identities}
\end{align}
The chart solves \(F_a(\Phi_a(z))=0\) and has retained
coordinates exactly \(z\).

Let \(G_R\) be the actual absolute contour gauge function.
It is not replaced by a background-relative penalty. Define
\begin{equation}
 \mathcal W_M(a)=S_M(e^{ia})+G_R(e^{ia}),\qquad
 f(a,z)=\frac1i\log[e^{i\Phi_a(z)}e^{ia}],
 \label{dmt-pulled-action}
\end{equation}
so \(f(a,0)=a\), and set
\begin{equation}
 \mathscr A_M(a)
    =D_z^2[\mathcal W_M(f(a,z))]_{z=0}.
 \label{dmt-independent-precision}
\end{equation}
This is the full analytic precision on the independent
background ball. Its Gaussian need not be centered at an
arbitrary independent background.

The actual physical branch \(a_*(w)\) satisfies
\[
 a_*(0)=0,\qquad F_R(e^{ia_*(w)})=e^{iw},\qquad
 P_R(e^{ia_*(w)};y,x)=I.
\]
It is the centered branch of the integral. Put
\begin{equation}
 A(w)=\mathscr A_M(a_*(w)),\qquad
 A_0=\mathscr A_M(0),\qquad C_0=A_0^{-1}.
 \label{dmt-physical-precision}
\end{equation}
The exact external logarithms are
\begin{equation}
 \Delta(w)=-\sum_c\log\det\pi_c(a_*(w)),\qquad
 \mathcal Q_{\det}(w)
       =-\frac12\log\det[A(w)A_0^{-1}],\qquad
 D_M=\Delta+\mathcal Q_{\det}.
 \label{dmt-leading-correction}
\end{equation}
All logarithmic branches in this display vanish at zero.
The factor \(R^{1-d}\) cancels between each pivot determinant
and its zero-background value.

The signs follow from the inverse pivot Jacobian and the
Gaussian factor \(\det A^{-1/2}\). There is no additional
background Haar factor \(j(a_*)\): the exact link Haar coordinates
are left-relative and give \(j(\Phi_a(g\xi))\), whose logarithm
vanishes at \(g=0\). The fluctuation straightening ratio is also
one at \(\xi=0\). Powers of \(g\), the gauge-insertion constant
and flat Gaussian constants are field independent and remain
outside \(D_M\). These definitions retain the actual mean,
composite map, chart and gauge; a direct-\(R\) replacement of
the constraint or an axial covariance would give a different
formula.

\subsection{The complete finite-regulator chain rule}
\label{dmt-chain-section}

For final coarse directions \(u,v\), write
\begin{equation}
 H_u=Da_*(0)u,\qquad B_{uv}=D^2a_*(0)[u,v].
 \label{dmt-stationary-responses}
\end{equation}
Define independent-background vertices by
\begin{align}
 \mathscr A_1[p]&=D\mathscr A_M(0)[p],&
 \mathscr A_2[p,q]&=D^2\mathscr A_M(0)[p,q],\notag\\
 \pi_{c,1}[p]&=D\pi_c(0)[p],&
 \pi_{c,2}[p,q]&=D^2\pi_c(0)[p,q].
 \label{dmt-background-vertices}
\end{align}
The physical derivatives are
\begin{align}
 A_u&=\mathscr A_1[H_u],&
 A_{uv}&=\mathscr A_2[H_u,H_v]+\mathscr A_1[B_{uv}],\notag\\
 \pi_{c,u}&=\pi_{c,1}[H_u],&
 \pi_{c,uv}&=\pi_{c,2}[H_u,H_v]+\pi_{c,1}[B_{uv}].
 \label{dmt-full-composition}
\end{align}

\begin{proposition}[Exact finite trace Hessian]
\label{dmt-finite-hessian}
For the actual finite integral,
\begin{align}
 D^2D_M(0)[u,v]
 ={}&-\sum_c\operatorname{tr}
       \{\pi_{c,uv}-\pi_{c,u}\pi_{c,v}\}\notag\\
 &-\frac12\operatorname{Tr}
       \{C_0A_{uv}-C_0A_uC_0A_v\}.
 \label{dmt-full-hessian}
\end{align}
Here lowercase trace acts on one pivot color block and
uppercase trace acts on all retained coordinates.
\end{proposition}
\begin{proof}
Differentiate the logarithmic determinants in
\eqref{dmt-leading-correction}, using
\[
 D\log\det A[B]=\operatorname{Tr}(A^{-1}B),
 \qquad D(A^{-1})[B]=-A^{-1}BA^{-1},
\]
and \(\pi_c(0)=I\). Products need not commute. Cyclicity of
the trace makes the answer symmetric in \(u,v\).
\end{proof}

In particular, constrained stationarity does not say that
the left current of \(S_M\) vanishes. It does not permit
deleting the constraint-curvature terms inside
\(\mathscr A_1\) or \(\mathscr A_2\).

\subsection{The second stationary response cancels after taking the scalar trace}
\label{dmt-color-section}

Introduce the independent-background scalar
\begin{equation}
 d_M(a)=-\sum_c\log\det\pi_c(a)
       -\frac12\log\det[\mathscr A_M(a)A_0^{-1}].
 \label{dmt-independent-scalar}
\end{equation}
This function is not asserted to be invariant under arbitrary
local background gauges. The constant-color symmetry suffices
for the next lemma.

\begin{lemma}[Semisimple cancellation of the first derivative]
\label{dmt-linear-cancellation}
For a semisimple Lie algebra,
\begin{equation}
                         Dd_M(0)=0.
 \label{dmt-zero-linear}
\end{equation}
Consequently the full Hessian in
\eqref{dmt-full-hessian} reduces to
\begin{align}
 D^2D_M(0)[u,v]
 ={}&-\sum_c\operatorname{tr}
 \{\pi_{c,2}[H_u,H_v]
       -\pi_{c,1}[H_u]\pi_{c,1}[H_v]\}\notag\\
 &-\frac12\operatorname{Tr}
 \{C_0\mathscr A_2[H_u,H_v]
       -C_0\mathscr A_1[H_u]C_0\mathscr A_1[H_v]\}.
 \label{dmt-reduced-hessian}
\end{align}
\end{lemma}
\begin{proof}
Let \(O=\operatorname{Ad}_k\) in the fixed real orthonormal
Lie basis. Endpoint covariance of the actual averages gives
\begin{equation}
                  \pi_c(Oa)=O\pi_c(a)O^{-1}.
 \label{dmt-pivot-color}
\end{equation}
The unique small pivot chart is equivariant because both
its constraint and retained-coordinate prescription are.
The absolute contour gauge is a sum of class functions,
and the actual scalar action is invariant under global
conjugation. Thus, with \(O_{\rm ret}\) the direct sum
on retained coordinates,
\begin{equation}
 \mathscr A_M(Oa)
       =O_{\rm ret}\mathscr A_M(a)O_{\rm ret}^{\,T},
 \qquad O_{\rm ret}^{\,T}O_{\rm ret}=I.
 \label{dmt-precision-color}
\end{equation}
The flat precision commutes with this representation.
Equations \eqref{dmt-pivot-color}--\eqref{dmt-precision-color}
prove invariance of \eqref{dmt-independent-scalar}, first
on the real domain and then by holomorphic continuation
with its zero-at-origin branches.

A semisimple Lie algebra has no invariant linear functional
on one adjoint color slot. Apply the invariance to a variation
on one current bond. Its linear functional is therefore zero.
Doing this for every bond proves \eqref{dmt-zero-linear}.
The argument applies separately to the pivot sum and the
Gaussian trace logarithm.

After substituting \eqref{dmt-full-composition} into
\eqref{dmt-full-hessian}, the coefficient of \(B_{uv}\)
is precisely \(Dd_M(0)[B_{uv}]\). This proves
\eqref{dmt-reduced-hessian}.
\end{proof}

This cancellation does not say that
\(\mathscr A_1=0\), \(\pi_{c,1}=0\), or that the current and
chart-curvature operators vanish. For an abelian factor,
retain \eqref{dmt-full-hessian} unless its linear functional
has been separately shown to vanish.

\subsection{The actual flat stationary lift}
\label{dmt-lift-section}

The first response appearing in
\eqref{dmt-reduced-hessian} is explicit in the proved flat
operators. Let
\begin{equation}
 C_R=D_a\{\log F_R(e^{ia})/i\}_{a=0},\qquad
 h=h_0,\quad K=K_0,\quad P_M=q_M+T_R^TT_R.
 \label{dmt-flat-data}
\end{equation}
Here \(q_M\) is the ordinary homogeneous classical Hessian,
and \(T_R\) is the linearized direct-\(R\) internal contour
map. The constraint in \(C_R\) is still the composite map.
The angle/counting normalization gives
\begin{equation}
 A_0=K^TP_MK,\qquad
 B_M=K^TP_Mh=K^Tq_Mh,\qquad
 H=h-KC_0B_M.
 \label{dmt-flat-lift}
\end{equation}
The simplification uses \(T_Rh=0\), since internal contours
use no pivots.

To verify the formula, use the identity-background affine
chart
\begin{align}
 \log F_R(e^{ia})/i&=C_Ra+\mathcal C_R(a),
 \qquad \mathcal C_R(0)=D\mathcal C_R(0)=0,\notag\\
 \mathfrak a(w,z)&=Kz+hw-h\mathcal C_R(\mathfrak a(w,z)).
 \label{dmt-affine-chart}
\end{align}
Put
\(\mathcal T_M(w,z)=\mathcal W_M(\mathfrak a(w,z))\).
Since the full action is stationary at zero,
\begin{equation}
 D_z\mathcal T_M(w,z)
       =A_0z+B_Mw+O((|w|+|z|)^2).
 \label{dmt-affine-stationarity}
\end{equation}
The positive flat retained Hessian and the established
local uniqueness identify its analytic critical point
with the physical branch. Thus
\(z_u=-C_0B_Mu\) and \(H_u=hu+Kz_u\), proving
\eqref{dmt-flat-lift}. No \(M^d\) or \(R^d\) transpose
factor belongs in these angle-coordinate matrices.

For completeness, the second response retained in the
general formula is also explicit. Set \(t_u=(u,z_u)\)
and define \(z_{uv}\) by the counting-pairing identity
\begin{equation}
 \langle\zeta,A_0z_{uv}\rangle
       =-D^3\mathcal T_M(0)[(0,\zeta),t_u,t_v]
       \quad\hbox{for every retained }\zeta.
 \label{dmt-second-stationary}
\end{equation}
Then
\begin{equation}
 B_{uv}=-hD^2\mathcal C_R(0)[H_u,H_v]+Kz_{uv}.
 \label{dmt-second-background}
\end{equation}
These equations follow by twice differentiating the retained
stationarity equation and \eqref{dmt-affine-chart}; they are
available even when the color cancellation is not.

\subsection{The complete precision vertices}
\label{dmt-precision-vertices-section}

Set
\begin{equation}
 E_a=D_zf(a,0),\qquad T_a=D_z^2f(a,0).
 \label{dmt-chart-tensors}
\end{equation}
With
\[
 J_aX=\int_0^1\operatorname{Ad}(e^{ita})X\,dt,\qquad
 \Xi_a(B)=\frac1i\log[e^{iB}e^{ia}],
\]
the exact tensors are
\begin{align}
 E_a&=J_a^{-1}K_a,\notag\\
 T_a[u,v]&=-J_a^{-1}h_aD_B^2F_a(0)[K_au,K_av]
       +D_B^2\Xi_a(0)[K_au,K_av].
 \label{dmt-both-chart-curvatures}
\end{align}
Consequently the arbitrary-background precision is
\begin{equation}
 \mathscr A_M(a)[u,v]
       =D^2\mathcal W_M(a)[E_au,E_av]
                  +D\mathcal W_M(a)[T_a(u,v)].
 \label{dmt-complete-precision}
\end{equation}
The two terms of \(T_a\) are respectively nonlinear constraint
curvature and conversion from the left-relative exponential
coordinate to the ordinary link logarithm. The absolute gauge
function remains inside \(\mathcal W_M\), so its nonphysical
background curvature is included as well. Its gradient
vanishes on the physical branch, but this does not permit
replacing the independent analytic family by a relative
quadratic gauge form.

Write \(W_n=D^n\mathcal W_M(0)\), so that
\(W_1=0\), \(W_2=P_M\), \(E_0=K\). Define
\[
 E_1[p]=D_aE_0[p],\qquad E_2[p,q]=D_a^2E_0[p,q],
 \qquad T_0=T_{a=0},\qquad T_1[p]=D_aT_0[p].
\]
For retained directions \(r,s\), direct differentiation gives
\begin{align}
 \mathscr A_1[p][r,s]
 ={}&W_3[p,Kr,Ks]
       +W_2[E_1[p]r,Ks]+W_2[Kr,E_1[p]s]\notag\\
 &+W_2[p,T_0(r,s)].
 \label{dmt-first-precision-vertex}
\end{align}
The last term is a first-background-order current and
constraint contribution; it remains despite \(W_1=0\).
The second vertex is
\begin{align}
 \mathscr A_2[p,q][r,s]
 ={}&W_4[p,q,Kr,Ks]\notag\\
 &+W_3[p,E_1[q]r,Ks]+W_3[p,Kr,E_1[q]s]\notag\\
 &+W_3[q,E_1[p]r,Ks]+W_3[q,Kr,E_1[p]s]\notag\\
 &+W_2[E_2[p,q]r,Ks]+W_2[Kr,E_2[p,q]s]\notag\\
 &+W_2[E_1[p]r,E_1[q]s]
       +W_2[E_1[q]r,E_1[p]s]\notag\\
 &+W_3[p,q,T_0(r,s)]\notag\\
 &+W_2[p,T_1[q](r,s)]
       +W_2[q,T_1[p](r,s)].
 \label{dmt-second-precision-vertex}
\end{align}
Only the \(W_1D_a^2T\) term disappears at zero.
The pivot vertices follow by twice differentiating
\begin{equation}
             \pi_c(a)=R^{d-1}(D_BF_a(0))_{c,p(c)}.
 \label{dmt-pivot-vertex-definition}
\end{equation}
All chart operations here use the actual finite-dimensional
logarithm and mean, and have finite range at the fixed \(R\).

\subsection{The classical cubic and quartic vertices}
\label{dmt-classical-vertices-section}

The classical vertices in
\eqref{dmt-first-precision-vertex}--\eqref{dmt-second-precision-vertex}
can be evaluated from the previously proved simultaneous
corrected-slice equations. This avoids an unproved bound for
the ordinary rooted average's physical adjoint.

Use the fine extension \(\mathcal H_M\), constrained Green
operator \(\mathcal C_M\), flat coarse Hessian \(q_M\),
corrected constraint remainder \(\mathcal R_M\), Wilson
jet \(f_\eta\), and its finite jet map \(\mathcal S\).
With the physical fine transposes and measures, the exact
pair \(Y=(A,m)\) obeys
\begin{align}
 A={}&\mathcal H_Ma-\mathcal H_M\mathcal R_M(A)
       -\mathcal C_M\mathcal S^Tf_\eta(\mathcal SA)
       +\mathcal C_MD\mathcal R_M(A)^Tm,\notag\\
 m={}&q_Ma-q_M\mathcal R_M(A)
       +(\mathcal S\mathcal H_M)^Tf_\eta(\mathcal SA)
       -\mathcal H_M^TD\mathcal R_M(A)^Tm.
 \label{dmt-classical-pair}
\end{align}
Write \(Y=\mathcal L_Ma+\mathcal V_M(Y)\), where
\(\mathcal L_M=(\mathcal H_M,q_M)\) and
\(\mathcal V_M(0)=D\mathcal V_M(0)=0\).
For \(\mathcal V_n=D^n\mathcal V_M(0)\), its responses are
\begin{align}
 Y_1[u]&=\mathcal L_Mu,\notag\\
 Y_2[u,v]&=\mathcal V_2[\mathcal L_Mu,\mathcal L_Mv],\notag\\
 Y_3[u,v,t]&=\mathcal V_3[
                  \mathcal L_Mu,\mathcal L_Mv,\mathcal L_Mt]\notag\\
 &\quad+\mathcal V_2[Y_2[u,v],\mathcal L_Mt]
       +\mathcal V_2[Y_2[u,t],\mathcal L_Mv]
       +\mathcal V_2[Y_2[v,t],\mathcal L_Mu].
 \label{dmt-classical-response-recursion}
\end{align}
These are derivatives, not Taylor coefficients, so no factorial
is omitted. The scalar envelope \(DS_M(a)=m(a)\) and equality
of the corrected and ordinary action values yield
\begin{align}
 D^3S_M(0)[u,v,t]
       &=\langle(Y_2[u,v])_m,t\rangle,\notag\\
 D^4S_M(0)[u,v,t,r]
       &=\langle(Y_3[u,v,t])_m,r\rangle.
 \label{dmt-classical-scalar-vertices}
\end{align}
Add the explicit local derivatives of \(G_R\) to obtain
\(W_3,W_4\).

These are finite compositions of the proved propagators and
finite-range analytic vertices. The uniform estimates use
complete physical cell spaces and smooth supported
reconstruction, not unweighted sums over \(M^4\) microscopic
coordinates. The fine \(\eta^4\) measure is already in the
physical transposes of \eqref{dmt-classical-pair}. It must
not be inserted again into the retained-coordinate traces.

\subsection{Absolute trace convergence and local kernels}
\label{dmt-convergence-section}

The strict precision theorem supplies symmetric pieces with
matrix rows and columns in their carrier and an analytic bound
of the following form, at fixed positive
\(\sigma,a,\rho\):
\begin{equation}
 \sup_i\sum_{Y\ni i}\sum_{j\in Y}
 e^{\sigma\tau(Y)+a|Y|}
 \|(\mathscr A_M-A_0)_{Y,ij}\|_\rho\le N.
 \label{dmt-strict-row-input}
\end{equation}
This is uniform in preceding depth and permitted volume.
The precision-composition theorem supplies the same form
after stationary substitution, keeping the actual matrix
rows rooted in their original carrier. The flat covariance
has an exponential Schur bound. The pivot scalar and
stationary field have their established strict and rooted
bounds.

\begin{lemma}[Source-rooted trace bounds]
\label{dmt-trace-convergence}
On a fixed smaller analytic domain, the actual Hessian
kernel
\[
 \mathcal K_{xy}=D^2D_M(0)[e_x,e_y]
\]
has an exponentially summable row bound
\begin{equation}
       \sup_x\sum_y e^{m d(x,y)}|\mathcal K_{xy}|
                       \le C_D
 \label{dmt-kernel-decay}
\end{equation}
for some \(m>0\). The traces in both
\eqref{dmt-full-hessian} and
\eqref{dmt-reduced-hessian} are absolutely convergent on
compact or summable variation directions. Constants may
depend on the fixed geometry, color and radii, but not
on preceding depth or volume.
\end{lemma}
\begin{proof}
A derivative in one background coordinate \(x\) acts only
on carriers containing \(x\). Cauchy bounds it by a fixed
constant times \(\rho^{-1}\) times the original majorant.
Two derivatives require both source slots in the carrier
and give \(\rho^{-2}\).

Reduce the spatial exponent in
\eqref{dmt-strict-row-input}. For each matrix row \(i\),
the unused exponent pays \(e^{-\delta d(i,x)}\), since
\(\tau(Y)\ge d(i,x)\) when \(i,x\in Y\). Its sum over
\(i\) is the finite lattice bound \(S_\delta\), also
uniformly on periodic graphs. With a second source \(y\),
spend another exponent portion on
\(\tau(Y)\ge d(x,y)\).

For the composed precision, this yields, after fixed
exponent reductions,
\begin{align}
 \sum_{ij}e^{m_1[d(i,x)+d(j,x)]}
                   \|(A_x)_{ij}\|&\le C_1,\notag\\
 \sum_{ij}\|(A_{xy})_{ij}\|
                   &\le C_2e^{-m_1d(x,y)}.
 \label{dmt-entry-source-bounds}
\end{align}
For the first line, use
\(d(i,x)+d(j,x)\le2\tau(Y)\) and reserve an additional
exponent for the row summation. Finite fiber and color
factors are absorbed into the constants.

The same argument before composition, followed by the
exponentially decaying external lift \(H\), bounds the
separate terms of \eqref{dmt-reduced-hessian}. Thus the
proof does not use cancellation of non-absolutely
convergent quantities.

The trace with \(A_{xy}\) is bounded by its absolute-entry
sum and the covariance entry envelope. For the other
trace expand
\begin{equation}
 \operatorname{Tr}(C_0A_xC_0A_y)
 =\sum_{i,j,k,l}\operatorname{tr}
     \{C_{0,ij}(A_x)_{jk}C_{0,kl}(A_y)_{li}\}.
 \label{dmt-trace-expansion}
\end{equation}
The source-root bounds and covariance envelope make this
absolutely convergent. The \(x\)-root weight at \(j\),
the covariance weight between \(j\) and \(i\), and the
\(y\)-root weight at \(i\) jointly pay \(d(x,y)\).
All remaining factors are absolutely summable. Hence
\[
 |\operatorname{Tr}(C_0A_xC_0A_y)|
                         \le C_3e^{-m_2d(x,y)}
\]
for a smaller positive \(m_2\). The pivot traces have the
same bound by local dependence before composition and the
rooted stationary family, or directly by their composed
strict scalar expansion. A final exponent reduction
allows the row sum in \eqref{dmt-kernel-decay}.
\end{proof}

Grouping retained coordinates into auxiliary \(J\)-cells
only regroups these traces; it inserts no extra factor
\(J^d\) into the exact identities.

On the infinite lattice a total log determinant at a
noncompact background need not exist as an extensive
number. The displayed formulas concern local kernels
and derivatives on compact or summable directions.
They are defined by the absolutely convergent sums above
or by a chosen joint local limit of the actual finite
matrices and vertices. The common majorants permit that
limit. No arbitrary finite activity is identified with
the periodization of an independently selected infinite
one.

\subsection{A convergent Maxwell moment of the traces}
\label{dmt-Maxwell-section}

Assume the chosen infinite total has the established
translation, hypercubic and local gauge identities, and
fix one simple color factor. The midpoint Hessian symbol
then has
\begin{equation}
 \widehat{\mathcal K}_{\mu\nu}^{ab}(p)
 =\delta^{ab}\beta(D_M)
       (|p|^2\delta_{\mu\nu}-p_\mu p_\nu)+O(|p|^4).
 \label{dmt-Maxwell-symbol}
\end{equation}
For \(\mu\ne\nu\), choose one trace-orthonormal color
vector \(t^a\), and let \(e_{x,\mu,a}\) be its final
unit-spacing link direction. The midpoint offsets
cancel for the diagonal orientation \(\mu,\mu\).
Therefore
\begin{equation}
 \boxed{\displaystyle
 \beta(D_M)=-\frac12\sum_{r\in\mathbb Z^4}r_\nu^2
       D^2D_M(0)[e_{0,\mu,a},e_{r,\mu,a}].}
 \label{dmt-Maxwell-moment}
\end{equation}
There is no color sum here. Hypercubic and color symmetry
make the choices equivalent. The minus sign comes from
twice differentiating \(e^{ip\cdot r}\); the factor
\(1/2\) extracts the coefficient of \(p_\nu^2\).
The actual classical reference has coefficient one in
this normalization.

Writing \(H_r=H e_{r,\mu,a}\), substitution of
\eqref{dmt-reduced-hessian} gives
\begin{align}
 \beta(D_M)
 ={}&\frac12\sum_r r_\nu^2\sum_c\operatorname{tr}
 \{\pi_{c,2}[H_0,H_r]
       -\pi_{c,1}[H_0]\pi_{c,1}[H_r]\}\notag\\
 &+\frac14\sum_r r_\nu^2\operatorname{Tr}
 \{C_0\mathscr A_2[H_0,H_r]
       -C_0\mathscr A_1[H_0]C_0\mathscr A_1[H_r]\}.
 \label{dmt-explicit-Maxwell-trace}
\end{align}
Every sum is absolutely convergent by
Lemma~\ref{dmt-trace-convergence}. Truncating the external
displacement at distance \(L\) has error at most
\(C_{m'}e^{-m'L}\) for any \(0<m'<m\).
Internal trace tails likewise have exponential bounds
after anchoring the sources and spending a spare exponent.

At an actual finite regulator, the trace Hessian formulas
remain exact. Apply the chosen symmetric short-lift
moment and divide by the moment of the actual finite
reference to obtain its finite Maxwell ratio. That
reference denominator is \(1+O(e^{-cN})\), not
automatically exactly one. No finite Ward identity at
discrete momenta is interpolated to obtain
\eqref{dmt-Maxwell-symbol}; its use requires the stated
infinite kernel and symmetries, or the separately
justified local-limit passage.

\subsection{What has and has not been reduced}
\label{dmt-scope-section}

The leading zero-input drift is expressed by actual local
pivot vertices, the flat stationary lift
\eqref{dmt-flat-lift}, the retained covariance, and the
full precision vertices
\eqref{dmt-first-precision-vertex}--\eqref{dmt-second-precision-vertex}.
The classical cubic and quartic vertices have the finite
operator recursion
\eqref{dmt-classical-pair}--\eqref{dmt-classical-scalar-vertices}.
All signs and coordinate normalizations are fixed.

The second physical stationary response was included
before being removed by the proved semisimple color
identity. The physical current was not replaced by an
independent parameter, and neither chart-curvature
contribution was omitted.

Evaluating \eqref{dmt-explicit-Maxwell-trace}, comparing
its values as \(M\) grows, determining a sign or a
universal coefficient, and controlling a physical
coupling trajectory remain further tasks. Absolute
convergence of the trace expression alone proves none
of those conclusions.
